\begin{exsol}{13}{Cabral 3.11 Smart car}{62}{Exam}{A2}
\asked Using \S3.4, identify the likely sources of strain in the jointly owned Daimler-Benz and
Swatch venture that produced the Smart car, which ended with Daimler buying out Swatch's share.

\method Firm boundaries as a balance between investment incentives (specific assets, hold-up)
and performance incentives; intermediate forms between market and hierarchy (Cabral \S3.4,
pp.~54--56). Supporting ideas: tacit knowledge and firm culture (\S3.5, p.~57); integration
trades hold-up for agency problems (\S13.1, pp.~338--339). Course contrast: NUMMI and
complementarities (\MR{} Ch.~4; notes, Aula~1, ``NUMMI'').

\solution
A joint venture is one of the intermediate forms of p.~55: neither a market contract nor a
single firm. \S3.4 predicts it inherits problems from both sides. The exercise gives only the
ownership history; the sources below are what the theory points to, not reported facts about
the partners' disputes.

\subparagraph{(1) Specific investments and hold-up between the owners.} Each partner brings
assets that, once committed, are valuable mainly inside the venture: a carmaker brings
automotive engineering and production, a watchmaker brings design and consumer-brand know-how.
Once sunk, each contribution is a specific asset,
and each partner can threaten to withhold cooperation to renegotiate the terms (the hold-up
problem of p.~55). A 50/50 ownership split also means each partner reaps only half the return on
any further investment it makes in the venture, so both under-invest relative to the joint
optimum. Item~30 (Cabral 13.10) computes the size of this effect: with two investors and an even
split, total surplus falls from $\tfrac12$ to $\tfrac38$.

\subparagraph{(2) Weak performance incentives inside, weak control outside.} A joint venture
removes the market discipline that a supplier relationship has (neither partner can easily
switch to another partner, p.~55) without creating the single line of authority of a firm.
With shared control, every major decision (design, cost targets, pricing, timing) must be
agreed by two owners whose objectives differ: each owner cares about the venture's profit and
also about its effect on its own core business and brand. Joint ownership with divergent
objectives is the agency problem of \S3.3 (p.~52) with two principals.

\subparagraph{(3) Two organizational systems that do not mix.} A carmaker and a watchmaker have
different capabilities, routines and cultures, much of them tacit (p.~57). A car designed around
one firm's ideas and built with the other's production system requires the two systems to be
merged inside a single product. The complementarity logic of Aula~1 says that a mix of practices
taken from two coherent systems can perform worse than either system alone, because each
practice is valuable only alongside its complements.

\subparagraph{(4) The resolution is the theory's fix.} Daimler's buyout ended joint control:
one owner now made all decisions and received the whole return on investment, the
integration solution to hold-up (p.~55; \S13.1, p.~338). The cost, as \S3.4 and \S13.1 warn, is
the loss of the partner's own incentive: Swatch's design and marketing input, if still used,
now had to be bought or managed from inside.

\subparagraph{(5) The contrast with NUMMI.} NUMMI was also a joint venture between carmakers
with different systems, but it was run on \emph{one} system: Toyota's production system governed
the plant, and GM's role was ownership and learning (notes, Aula~1). The internal strain
predicted by points~(2)--(3) was therefore small, since there was no merging of two systems
inside the operation. NUMMI's limitation was elsewhere: GM could not transfer the system
piece by piece to its own plants, which is again the complementarity point. Smart illustrates
the opposite configuration: the systems had to be merged within the venture itself, so the
strain appeared inside it.

\interp Joint ventures look like a way to combine the best of two firms, but \S3.4 says they
combine the hold-up exposure of separate ownership with the weak incentives of shared control.
They work best when each partner's contribution is separable or when one partner's system
clearly governs operations (NUMMI); they strain when two systems must be fused (Smart), and
the usual exit is that one partner buys the other out.

\verify Cross-check against the theory's predictions. If strain came from hold-up and shared
control, then (i) the venture should end with consolidation of ownership rather than with a
looser market contract, and (ii) the natural buyer is the partner in whose core business the
product lies (a car lies in the carmaker's). Both match the outcome reported in the exercise
(Daimler took complete control in November 1998).

\answer{\S3.4 points to (1) specific investments by both owners and hold-up between them, with a
50/50 split under-rewarding each partner's investment; (2) shared control between owners with
different objectives, i.e.\ weak incentives without a single authority; and (3) two different,
largely tacit organizational systems that had to be fused in one product. Daimler's buyout is
the integration fix. NUMMI, run on Toyota's single system, avoided the internal strain.}
\end{exsol}
